\chapter{Introduction}
\label{chap:introduction}

\section{Background}
\label{section:background}

Character creation has become a fundamental aspect of the modern gaming experience, particularly in Role-Playing Games (RPGs) and multiplayer titles. As hardware capabilities have advanced, character creation systems have grown increasingly detailed, offering players greater freedom to express their identity and personal aesthetic through their in-game avatar.

Unreal Engine, developed by Epic Games, is one of the most widely adopted game development platforms in the industry, used across AAA productions, independent studios, and academic projects alike. Its powerful rendering capabilities and support for high-fidelity character systems such as MetaHuman has made it a popular choice for games that emphasize detailed character representation. Unreal Engine also supports a plugin architecture that allows developers to extend its functionality without modifying the engine's core codebase, making it a practical platform for developing reusable, drop-in tools.

Currently, most video game character creation systems rely on a slider and dropdown paradigm, where players manually adjust individual appearance parameters such as jaw width, cheek height, eye spacing, and skin tone. While this approach gives players precise control, the complexity scales directly with the number of available parameters. As character systems become more detailed, players are presented with an ever-growing array of sliders and options.

Recent advances in Large Language Models (LLMs) have demonstrated that natural language can serve as a reliable interface for producing structured, machine-readable output. Models such as Llama, GPT-4o, and Gemini are capable of interpreting open-ended descriptive text and returning well-formed JSON with consistent values — a capability that has matured significantly in recent years through techniques such as constrained decoding and few-shot prompting.

\section{Problem Statement}
\label{section:problem-statement}

As established in the background, current character creation systems can become overwhelming and unintuitive when players are faced with a large number of sliders and abstract numerical values. This complexity often causes confusion and requires significant time investment before a player can produce a character that matches their expectations. Furthermore, players who lack familiarity with facial anatomy or artistic terminology may struggle to understand what parameters like "zygomatic width" or "nasal bridge depth" actually affect visually, frequently resulting in a character that does not reflect their original intent.

While reusable character creation systems such as Unreal Engine's Mutable and MetaHuman already exist for game developers, none of them provide a natural language interface. There is currently no standardized, reusable tool within the Unreal Engine ecosystem that bridges natural language input with existing character parameter systems.

\section{Solution Overview}
\label{section:solution-overview}

To address these pain points, this project proposes an Unreal Engine plugin that bridges a player's natural language description and the underlying character parameter system. When a player types a description such as "a tall warrior with a broad jaw and weathered skin," the plugin passes that input to a Large Language Model, which interprets it and returns a structured set of character parameter values that are then automatically applied to the character mesh.

The plugin is designed as a drop-in tool for game developers. Once integrated, it handles all LLM communication and response parsing internally, supporting both locally running models for offline use and cloud-based models for higher output quality, with automatic fallback between the two.

\newpage
\subsection{Features}
\label{subsection:features}

\subsubsection{For Game Developers (directly)}
\begin{enumerate}
    \item \textbf{Drop-in Plugin Integration}: The plugin can be added to any Unreal Engine project without modifying the engine or existing game code.
    
    \item \textbf{Parameter Registration}: Developers register their character parameters and map them to their character mesh through a configuration asset.
    
    \item \textbf{Preset Configurations}: Built-in presets for popular character systems are provided to reduce initial setup time.
    
    \item \textbf{LLM Mode Selection}: Developers choose between local inference or cloud-based API at setup time based on their project requirements.
    
    \item \textbf{Automatic Fallback}: The plugin automatically falls back to alternative LLM options or default values if the primary mode fails.
    
    \item \textbf{Blueprint Event Integration}: The plugin fires a Blueprint event with parsed parameter values, allowing developers to handle character application in their own way.
\end{enumerate}

\subsubsection{For Players (indirectly)}
\begin{enumerate}
    \item \textbf{Natural Language Input}: Players describe their desired character appearance in plain text instead of manually adjusting sliders and numerical values.
    
    \item \textbf{Automatic Character Generation}: Character appearance parameters are interpreted and applied automatically from the player's text description.
    
    \item \textbf{Flexible Description Style}: The plugin accepts both explicit descriptions and implied or nuanced expressions.
\end{enumerate}

\newpage
\section{Target User}
\label{section:target-user}

\textbf{Game Developer (Direct User)}
\begin{itemize}
    \item \textbf{Demographic}: Independent developers or small 
    studio teams building games with Unreal Engine across any genre 
    that features character creation.
    
    \item \textbf{Technical Skill}: Expected to have working 
    familiarity with Unreal Engine's Blueprint or C++ workflow. 
    Responsible for integrating and configuring the plugin within 
    their project, including parameter registration and mapping.
    
    \item \textbf{Industry or Domain}: Game development industry, 
    specifically projects built on Unreal Engine that require a 
    character creation system.
\end{itemize}

\textbf{Player (Indirect User)}
\begin{itemize}
    \item \textbf{Demographic}: End users who interact with the 
    character creation screen at the start of the game. May range 
    from casual players to experienced gamers across any age group.
    
    \item \textbf{Technical Skill}: No technical knowledge required. 
    Players interact with the plugin's output without any awareness 
    of its underlying implementation.
    
    \item \textbf{Industry or Domain}: General gaming audience, 
    particularly players of RPGs or any game genre that features 
    character customization.
\end{itemize}

\section{Terminology}
\label{section:terminology}

\begin{description}
    \item[\textbf{Large Language Model (LLM)}] A type of artificial 
    intelligence model trained on large amounts of text data, capable 
    of understanding and generating human language. In this project, 
    LLMs are used to interpret player descriptions and produce 
    structured character parameter values.

    \item[\textbf{Morph Target / Blend Shape}] A 3D animation 
    technique that stores a deformed version of a mesh. By blending 
    between the original and deformed states using a numerical value 
    between 0 and 1, specific facial or body features such as jaw 
    width or cheek height can be adjusted at runtime.

    \item[\textbf{Structured Output}] Output generated by an LLM 
    that conforms to a predefined format, typically JSON. Structured 
    output ensures that the model's response can be reliably parsed 
    and used by the application without manual interpretation.

    \item[\textbf{Prompt Engineering}] The practice of designing and 
    refining the input given to an LLM in order to guide it toward 
    producing accurate and consistent outputs. This project uses 
    prompt engineering as the primary method of controlling LLM 
    behavior without requiring model training.

    \item[\textbf{Few-shot Prompting}] A prompt engineering technique 
    where a small number of input-output examples are included in the 
    prompt to demonstrate the expected behavior to the model. This 
    allows the model to generalize to new inputs without any 
    additional training.

    \item[\textbf{Constrained Decoding}] A technique used in local 
    LLM inference that restricts the model's output to conform to a 
    defined grammar or format, such as a GBNF grammar file. This 
    guarantees well-formed structured output even from smaller local 
    models.

    \item[\textbf{Unreal Engine Plugin}] A self-contained module that 
    extends Unreal Engine's functionality without modifying the 
    engine's core codebase. Plugins can be dropped into any 
    compatible Unreal Engine project and activated through the 
    engine's plugin manager.

    \item[\textbf{Blueprint}] Unreal Engine's visual scripting system 
    that allows developers to implement game logic without writing 
    code directly. This project exposes plugin functionality through 
    Blueprint events so that developers can integrate it regardless 
    of their preferred workflow.

    \item[\textbf{Skeletal Mesh}] A 3D character mesh bound to a 
    skeleton, allowing it to be animated and deformed. Morph targets 
    are applied on top of a skeletal mesh to modify the character's 
    appearance.

    \item[\textbf{Data Asset}] An Unreal Engine asset type used to 
    store structured configuration data. This project uses a Data 
    Asset to store the developer's parameter registration and blend 
    shape mapping configuration.

    \item[\textbf{Semantic Parameter}] A human-readable parameter 
    name that describes a character feature conceptually, such as 
    \texttt{jaw\_width} or \texttt{eye\_spacing}. Semantic parameters 
    are defined by the developer and mapped to the actual blend shape 
    names used by their character mesh.

    \item[\textbf{Local Inference}] Running an LLM directly on the 
    player's machine without requiring an internet connection. This 
    project supports local inference via a locally hosted model 
    server, allowing the plugin to function in offline environments.

    \item[\textbf{Cloud API}] A remote service accessed over the 
    internet that provides LLM capabilities. Cloud APIs typically 
    offer higher output quality than local models at the cost of 
    requiring an internet connection and an API key provided by the 
    developer.
\end{description}